\begin{tikzpicture}[
  font=\sffamily,
  card/.style={
    rectangle, rounded corners=4pt, draw=ceindigo, line width=0.9pt,
    fill=cepanel, text width=4.3cm, align=center, inner sep=7pt, minimum height=1.85cm
  },
  comp/.style={
    rectangle, rounded corners=3pt, draw=cemuted, line width=0.7pt,
    fill=cecream, text width=5.95cm, align=left, inner sep=7pt
  }
]
  % Fila superior: Comparación Causal vs Efectual
  \node[comp, draw=ceslate] (caus) at (-0.25, 3.3) {
    \textcolor{ceslate}{\bfseries LÓGICA CAUSAL (Planificación clásica)}\\[3pt]
    \scriptsize Parte de un \textbf{objetivo fijo} predeterminado y selecciona los medios necesarios para alcanzarlo.\\
    $\bullet$ Predicción del futuro mediante análisis previo.\\
    $\bullet$ Óptima en mercados maduros y estables.
  };
  \node[comp, draw=ceteal, fill=cepanelteal] (efec) at (7.35, 3.3) {
    \textcolor{ceteal!90!black}{\bfseries LÓGICA EFECTUAL (Sarasvathy, 2001)}\\[3pt]
    \scriptsize Parte de los \textbf{medios disponibles hoy} y explora qué objetivos emergen al actuar con otros.\\
    $\bullet$ Control del futuro mediante acción iterativa.\\
    $\bullet$ Óptima bajo incertidumbre radical (startups).
  };

  \draw[<->, >=Stealth, thick, ceamber] (caus.east) -- (efec.west) node[midway, fill=white, inner sep=2pt, font=\sffamily\bfseries\scriptsize, ceamber!90!black] {vs.};

  % Los 5 principios de la Efectuación en abanico inferior
  \node[font=\sffamily\bfseries\small, ceindigo] at (3.55, 1.75) {LOS CINCO PRINCIPIOS DE LA EFECTUACIÓN};

  \node[card] (p1) at (-1.5, 0.25) {
    \textcolor{cecrimson}{\bfseries\footnotesize 1. PÁJARO EN MANO}\\[2pt]
    \scriptsize Partir de quién soy, qué sé y a quién conozco, sin esperar al capital ideal.
  };
  \node[card] (p2) at (3.55, 0.25) {
    \textcolor{cecrimson}{\bfseries\footnotesize 2. PÉRDIDA ACEPTABLE}\\[2pt]
    \scriptsize Acotar cuánto puedo arriesgar en cada paso antes que estimar retornos lejanos.
  };
  \node[card] (p3) at (8.6, 0.25) {
    \textcolor{cecrimson}{\bfseries\footnotesize 3. COLCHA DE RETAZOS}\\[2pt]
    \scriptsize Tejer alianzas tempranas con socios comprometidos que co-crean el proyecto.
  };

  \node[card, fill=cepanelamber, draw=ceamber] (p4) at (1.0, -2.0) {
    \textcolor{ceamber!90!black}{\bfseries\footnotesize 4. PRINCIPIO DE LIMONADA}\\[2pt]
    \scriptsize Convertir sorpresas, errores y contingencias en palancas de nuevo valor.
  };
  \node[card, fill=cepanelamber, draw=ceamber] (p5) at (6.1, -2.0) {
    \textcolor{ceamber!90!black}{\bfseries\footnotesize 5. PILOTO EN EL AVIÓN}\\[2pt]
    \scriptsize El futuro no se predice con estadísticas: se fabrica mediante decisiones propias.
  };
\end{tikzpicture}
